\documentclass[conference]{IEEEtran}

% ============================================================================
% PACKAGES
% ============================================================================
\usepackage{cite}
\usepackage{amsmath,amssymb,amsfonts}
\usepackage{algorithmic}
\usepackage{graphicx}
\usepackage{textcomp}
\usepackage{xcolor}
\usepackage{booktabs} % For professional looking tables
\usepackage{multirow}
\usepackage{float}    % To force image placement

\def\BibTeX{{\rm B\kern-.05em{\\sc i\kern-.025em b}\kern-.08em
    T\kern-.1667em\lower.7ex\hbox{E}\kern-.125emX}}

\begin{document}

% ============================================================================
% TITLE & AUTHORS
% ============================================================================
\title{Deep Learning-Based Reconstruction and Analysis of ANSYS Finite Element Bone Meshes Using PointNet Architecture}

% TODO: FILL IN YOUR ACTUAL NAME AND SUPERVISOR NAME HERE
% \author{
% \IEEEauthorblockN{Your Name}
% \IEEEauthorblockA{\textit{Department of [Your Department]} \\
% \textit{[Your University]}\\
% [City, Country] \\
% your.email@university.edu}
% \and
% \IEEEauthorblockN{Dr. Mahmuda [Supervisor Last Name]}
% \IEEEauthorblockA{\textit{Department of [Supervisor's Department]} \\
% \textit{[Your University]}\\
% [City, Country] \\
% supervisor.email@university.edu}
% }

\maketitle

% ============================================================================
% ABSTRACT
% ============================================================================
\begin{abstract}
Finite Element Analysis (FEA) is the gold standard for biomechanical simulation but suffers from high computational costs and manual preprocessing bottlenecks. This paper presents a novel framework bridging the gap between industrial simulation software (ANSYS) and Geometric Deep Learning. We introduce an automated end-to-end pipeline that parses raw ANSYS \texttt{.cdb} files, extracts unstructured nodal data, and utilizes a PointNet AutoEncoder for mesh quality assessment and classification. Unlike standard 3D vision tasks, our approach validates reconstruction not just visually, but through engineering metrics including volume preservation and Hausdorff distance. Experimental results demonstrate a Chamfer distance of 0.201 and 84.2\% volume preservation, proving that deep learning can effectively act as a surrogate model for rapid bone mesh analysis.
\end{abstract}

\begin{IEEEkeywords}
PointNet, ANSYS, Finite Element Analysis, Geometric Deep Learning, Biomechanics.
\end{IEEEkeywords}

% ============================================================================
% I. INTRODUCTION
% ============================================================================
\section{Introduction}
Finite Element Analysis (FEA) is critical in computational biomechanics for evaluating bone stress distributions, fracture risks, and implant stability. However, the traditional FEA workflow—ranging from mesh generation to solver execution—is computationally expensive and relies heavily on manual expert intervention for quality assessment. As the volume of patient-specific medical imaging data increases, the inability to rapidly process and analyze these simulations becomes a critical bottleneck.

Deep Learning (DL), specifically Geometric Deep Learning (GDL), offers a promising solution by learning directly from unstructured 3D data. While Convolutional Neural Networks (CNNs) excel at regular grids (images), they fail to efficiently capture the sparse, irregular topology of Finite Element meshes. 

This research proposes an automated framework that bridges the gap between industrial simulation software (ANSYS) and modern Deep Learning. We introduce a pipeline that parses raw ANSYS APDL command data, extracts nodal coordinates, and utilizes a PointNet-based architecture to learn latent geometric representations of bone structures. This approach enables rapid, automated quality assessment and classification, laying the groundwork for real-time surrogate modeling in clinical environments.

% ============================================================================
% II. BACKGROUND STUDY
% ============================================================================
\section{Background Study}

\subsection{FEA in Biomechanics}
Finite Element models are extensively used to study trabecular and cortical bone mechanics \cite{fea_biomech}. The standard workflow involves utilizing software like ANSYS or Abaqus to generate volumetric meshes. However, analyzing the quality of these meshes (e.g., checking for distorted elements or volume preservation) is traditionally a post-process manual step that requires significant domain expertise.

\subsection{Deep Learning on 3D Point Clouds}
Early approaches to 3D learning involved voxelization (VoxNet \cite{voxnet}), which converts meshes into 3D pixel grids. This method suffers from high memory consumption $\mathcal{O}(N^3)$ and loss of fine geometric detail crucial for medical analysis. 

A breakthrough occurred with PointNet \cite{pointnet}, which consumes raw point clouds directly. By utilizing symmetric functions (max pooling) and Multi-Layer Perceptrons (MLPs) shared across points, PointNet respects the permutation invariance of 3D data. While PointNet has been widely applied to standard CAD datasets like ModelNet40 or ShapeNet \cite{shapenet}, its application to raw, unstructured industrial FEA file formats remains underexplored.

% ============================================================================
% III. RESEARCH GAP
% ============================================================================
\section{Research Gap}
Despite advancements in both FEA and Computer Vision, a significant disconnect remains:
\begin{enumerate}
    \item \textbf{Data Incompatibility:} Current Geometric DL models operate on clean, pre-processed 3D formats (e.g., .obj, .ply). There is a lack of frameworks capable of directly ingesting raw industrial simulation files (ANSYS \texttt{.cdb}) which contain simulation metadata mixed with geometry.
    \item \textbf{Lack of FEA-Specific Metrics:} Most 3D reconstruction research focuses on visual fidelity. There is a gap in evaluating reconstruction models based on engineering metrics relevant to biomechanics, such as Volume Preservation and physical node fidelity.
    \item \textbf{Manual Bottlenecks:} The assessment of bone mesh quality is currently not automated within the Deep Learning loop, preventing the creation of large-scale, self-optimizing biomechanical databases.
\end{enumerate}

This thesis addresses these gaps by developing a custom parser-to-prediction pipeline specifically for ANSYS-generated bone meshes.

% ============================================================================
% IV. PROBLEM DOMAIN
% ============================================================================
\section{Problem Domain}
This research focuses on the intersection of \textbf{Computational Biomechanics} and \textbf{Unsupervised Representation Learning}. The specific scope includes:
\begin{itemize}
    \item \textbf{Input Domain:} Raw ANSYS \texttt{.cdb} (Command Database) files containing \texttt{NBLOCK} (Node Block) data structures representing human bone geometries.
    \item \textbf{Task 1: Latent Representation \& Reconstruction:} Utilizing a PointNet AutoEncoder to compress the high-dimensional mesh data into a compact latent vector and reconstruct it to verify that the model captures the complex topology of bone structures.
    \item \textbf{Task 2: Classification:} Automated identification of bone side (Left/Right) and structural categories directly from the point cloud data without manual labeling.
    \item \textbf{Output Metrics:} Evaluation is based on Chamfer Distance (geometric accuracy), Hausdorff Distance (outlier handling), and Volume Preservation (physical fidelity).
\end{itemize}

% ============================================================================
% V. METHODOLOGY
% ============================================================================
\section{Methodology}

\subsection{Data Parsing Pipeline (Novelty)}
The core contribution of this work is the custom parsing engine designed to interpret ANSYS APDL formats. Unlike standard 3D loaders, our system reads the \texttt{NBLOCK} format:
\begin{equation}
    P = \{ (x_i, y_i, z_i) \mid i \in 1...N \}
\end{equation}
where $P$ represents the extracted point cloud. The parser filters simulation commands (EBLOCK, MAT) to isolate pure geometric nodes, converting them into a normalized tensor format suitable for Deep Learning.

\subsection{Network Architecture}
We utilize a modified PointNet AutoEncoder. The encoder consists of purely convolutional layers (64, 128, 1024 filters) followed by a Global Max Pooling layer to extract the global feature vector $G$. The decoder attempts to reconstruct the input point cloud from $G$ using fully connected layers (1024, 2048, $N \times 3$).

% PLACEHOLDER FOR DIAGRAM
% \begin{figure}[htbp]
% \centering
% \includegraphics[width=0.9\linewidth]{your_pipeline_image.png}
% \caption{Proposed End-to-End Pipeline: From ANSYS CDB to PointNet Analysis.}
% \label{fig:pipeline}
% \end{figure}

\subsection{Loss Function}
To ensure the reconstructed mesh aligns with the physical reality of the bone, we minimize the Chamfer Distance (CD):
\begin{equation}
    CD(S_1, S_2) = \sum_{x \in S_1} \min_{y \in S_2} ||x-y||_2^2 + \sum_{y \in S_2} \min_{x \in S_1} ||x-y||_2^2
\end{equation}
This metric penalizes points that drift from the original FEA node locations.

% ============================================================================
% VI. EXPERIMENTAL RESULTS
% ============================================================================
\section{Experimental Results}

\subsection{Quantitative Analysis}
The model was trained on a dataset of patient-specific bone meshes. Table I summarizes the reconstruction performance. The low Chamfer Distance (0.2009) indicates high geometric fidelity, while the F1 Score (0.7230) demonstrates the model's ability to retain key structural features.

\begin{table}[htbp]
\caption{Reconstruction Performance on Bone Mesh Dataset}
\begin{center}
\begin{tabular}{lcc}
\toprule
\textbf{Metric} & \textbf{Mean} & \textbf{Std Dev} \\
\midrule
Chamfer Distance & 0.2009 & 0.0062 \\
Hausdorff Distance & 0.0430 & 0.7556 \\
F1 Score & 0.7230 & 0.0216 \\
Volume Preservation & 84.2\% & 5.8\% \\
\bottomrule
\end{tabular}
\label{tab1}
\end{center}
\end{table}

\subsection{Qualitative Analysis}
Visual inspection confirms that the PointNet AutoEncoder successfully captures the complex curvature of the bone epiphysis and diaphysis.

% PLACEHOLDER FOR RESULT IMAGE
% \begin{figure}[htbp]
% \centering
% \includegraphics[width=0.8\linewidth]{your_result_image.png}
% \caption{Comparison of Original ANSYS Mesh (Left) vs. Deep Learning Reconstruction (Right).}
% \label{fig:results}
% \end{figure}

% ============================================================================
% VII. CONCLUSION
% ============================================================================
\section{Conclusion}
This study presented a framework for integrating ANSYS finite element data with geometric deep learning. By successfully parsing and reconstructing industrial bone meshes with high fidelity (84.2\% volume preservation), we demonstrated that Deep Learning can effectively serve as a surrogate model for biomechanical analysis. Future work will focus on predicting stress tensors directly from the point cloud, further reducing the computational cost of FEA simulations.

% ============================================================================
% REFERENCES
% ============================================================================
\begin{thebibliography}{00}

\bibitem{fea_biomech} U. Wolfram and H. Gross, ``Finite Element Analysis in Biomechanics,'' in \textit{Computational Modeling in Biomechanics}, Springer, 2010.

\bibitem{pointnet} C. R. Qi, H. Su, K. Mo, and L. J. Guibas, ``PointNet: Deep Learning on Point Sets for 3D Classification and Segmentation,'' in \textit{Proc. IEEE CVPR}, 2017.

\bibitem{voxnet} D. Maturana and S. Scherer, ``VoxNet: A 3D Convolutional Neural Network for Real-Time Object Recognition,'' in \textit{Proc. IEEE/RSJ IROS}, 2015.

\bibitem{shapenet} Z. Wu et al., ``3D ShapeNets: A Deep Learning Approach for 3D Shape Representation,'' in \textit{Proc. IEEE CVPR}, 2015.

\end{thebibliography}

\end{document}